\documentclass[12pt]{article}
\usepackage[margin=1in]{geometry}

\usepackage{hyperref}
\usepackage{listings}
\usepackage{rotating, graphicx}
\usepackage{booktabs, natbib}

\usepackage{amsmath, enumerate, quoting, amssymb, bm}
%\usepackage{configuration, subfigure, url}
\usepackage{subfigure, url}

\usepackage{color}

\definecolor{blue2}{rgb}{0.1, 0.2, 0.6}
%
%\renewcommand{\red}[1]{\textcolor{red}{#1}}
%\renewcommand{\blue}[1]{\textcolor{blue}{#1}}
\newcommand{\gv}[1]{\textcolor{blue}{(GV: #1)}}
\newcommand{\jy}[1]{\textcolor{red}{(JY: #1)}}
\newcommand{\why}[1]{\textcolor{red}{(WHY: #1)}}
\newcommand{\zhy}[1]{\textcolor{green}{(ZHY: #1)}}
\DeclareMathOperator*{\argmin}{arg\,min}
\DeclareMathOperator*{\cov}{Cov}
\DeclareMathOperator*{\corr}{Corr}

\hypersetup{
	colorlinks=true,
	linkcolor=blue,
	filecolor=blue,
	citecolor = blue,
	urlcolor=cyan,
}

%\usepackage{xr}
%\externaldocument{isgee}
%\externaldocument{supp}

% remove left indentation in itemize
\usepackage{enumitem}
\setlist[itemize]{leftmargin=*}

% graph path
\graphicspath{{./figs/}}

\quotingsetup{leftmargin = 0pt}

\newenvironment{comment}%
{\begin{quoting}\noindent\small\it\color{blue2}\ignorespaces%
  }{\end{quoting}}




\begin{document}

\begin{center}
  {\Large\bf Response to the Comments}
\end{center}

\section*{Summary}

We thank the Editor and the referee for the time spent reviewing our manuscript 
and for the insightful comment. We have revised the paper to address this 
comment and hope that the changes meet your expectations. We look forward 
to the possibility of our article appearing in your journal.

\section*{Reviewer 1}

\begin{comment}
Specific comment: In the developed model, the ‘baseline rate of occurrence’ is 
assumed to be constant $r$. Can the authors extend their model to the case when 
it is a function of time, say, $r(t)$? This extension will also have meaningful 
practical meaning as the division rate of a stem cell may change in time as it 
ages? I assume that, under this generalized setting, the proof of Proposition 1 
can be derived similarly?
\end{comment}

The referee is absolutely correct. The proof of Proposition 1 continues to hold 
even when $r$ is a deterministic function of time, provided it satisfies 
appropriate regularity conditions that allow the use of the Dominated Convergence 
Theorem. We have added a remark immediately following the proposition that includes 
modified formulas for both the expectation and the variance of the stem cell count. 
In the same remark, we also discuss why, for our particular data sets, 
the assumption of a constant division rate remains a practical and reasonable 
modeling choice.

\end{document}

%%% Local Variables:
%%% mode: latex
%%% TeX-master: t
%%% End:
